% AUTO-GENERATED by scripts/exp_sync_kernel_output_potential_branch_radius_certificate.py
\begin{proposition}[复参数分歧点的判别式刻画]\label{prop:pressure-branchpoints-discriminant}
把 $F(\lambda,u)\in\ZZ[u][\lambda]$ 视为关于 $\lambda$ 的代数方程，并在 $u\in\CC^\times$ 上延拓其代数函数分支。
则 $\lambda(u)$ 的分歧点恰对应 $\lambda$ 变为重根，即存在 $(\lambda,u)\in\CC\times\CC^\times$ 使得
$$
F(\lambda,u)=0,\qquad \partial_{\lambda}F(\lambda,u)=0.
$$
等价地，分歧点的 $u$-集合由关于 $\lambda$ 的判别式零点给出；对本附录的六次 $F$，有完全显式分解
$$
\mathrm{Disc}_{\lambda}(F)(u)=256 u^{25} + 1408 u^{24} + 22144 u^{23} - 20217 u^{22} + 814606 u^{21} + 5139405 u^{20} + 10776760 u^{19} + 16410991 u^{18} + 23046562 u^{17} + 30590013 u^{16} + 34775088 u^{15} + 30590013 u^{14} + 23046562 u^{13} + 16410991 u^{12} + 10776760 u^{11} + 5139405 u^{10} + 814606 u^{9} - 20217 u^{8} + 22144 u^{7} + 1408 u^{6} + 256 u^{5}=-u^{5}\,D(u),
$$
其中 $D(u)\in\ZZ[u]$ 为 $20$ 次回文多项式：
$$
D(u)=256 u^{20} + 1408 u^{19} + 22144 u^{18} - 20217 u^{17} + 814606 u^{16} + 5139405 u^{15} + 10776760 u^{14} + 16410991 u^{13} + 23046562 u^{12} + 30590013 u^{11} + 34775088 u^{10} + 30590013 u^{9} + 23046562 u^{8} + 16410991 u^{7} + 10776760 u^{6} + 5139405 u^{5} + 814606 u^{4} - 20217 u^{3} + 22144 u^{2} + 1408 u + 256.
$$
\end{proposition}

\begin{corollary}[最近分歧点与 $\theta$-解析半径]\label{cor:pressure-analytic-radius}
以 $\theta=\log u$ 为局部坐标，并以 $\lambda(1)=3$ 的 Perron 分支为基准进行解析延拓。
令 $u_\star$ 为 $D(u)=0$ 的根中、使得 $|\log u|$（在 $\log u$ 的 $2\pi i\ZZ$ 选取中取最小模）最小者。
数值上可取
$$
u_\star\approx 0.5898937925+1.0420567812i,\qquad u_\star^{-1}\approx 0.4114034967-0.7267508304i,
$$
并且相应的最近分歧点可取
$$
\theta_\star=\log u_\star\approx 0.1801840156+1.0556869306i,\qquad -\theta_\star,
$$
从而 $\theta=0$ 处解析芽的最大圆盘半径为
$$
\boxed{\ R_\theta:=|\theta_\star|\approx 1.0709533953\ }.
$$
因此在 $|\theta|<R_\theta$ 内，$\lambda(e^{\theta})$ 与 $P(\theta)=\log\lambda(e^{\theta})$ 可作为单值解析函数延拓；并且在 $|\theta|=R_\theta$ 处发生代数分歧（分支点）。
\end{corollary}

\begin{corollary}[Taylor 截断余项的 Cauchy 证书]\label{cor:pressure-taylor-remainder-cauchy}
设 $T_8(\theta)$ 为 $P(\theta)$ 在 $\theta=0$ 的 Taylor 多项式截断到 $\theta^8$。取任意 $0<r<R_\theta$ 并记
$$
M_r:=\max_{|\theta|=r}|P(\theta)|.
$$
则对任意 $|\theta|<r$ 有一致余项上界
$$
\boxed{\ \bigl|P(\theta)-T_8(\theta)\bigr|\le M_r\cdot \frac{|\theta|^{9}}{r^{9}}\cdot\frac{1}{1-|\theta|/r}\ }.
$$
取 $r:=0.99\,R_\theta\approx 1.060244$ 并沿 $|\theta|=r$ 对分支作连续数值跟踪（初值 $\lambda(1)=3$），得到上界 $M_r\le 1.707619$。
因此当 $|\theta|\le 0.5$ 时，有可审计余项界
$$
\bigl|P(\theta)-T_8(\theta)\bigr|\le 3.728189e-03.
$$
\end{corollary}

\begin{remark}[分歧半径的势指纹：与负携带势的对比]\label{rem:pressure-radius-compare-phi-minus}
在附录 \ref{app:vector-potential} 的负携带势实例中，主特征值满足三次方程
$$
\lambda^3-(u+2)\lambda^2+(u-2)\lambda+3u=0,\qquad u=e^{\theta}.
$$
对该三次关于 $\lambda$ 取判别式（忽略非零常数因子）得到分歧点集合由四次多项式
$$
13 u^{4} + 14 u^{3} - 83 u^{2} + 264 u + 48=0
$$
刻画；其最近分歧点可取
$$
u_\star^{(-)}\approx 1.5189501269-1.7690244752i,\qquad \theta_\star^{(-)}\approx 0.8465821825-0.8613092768i,
$$
从而对应解析半径
$$
R_\theta^{(-)}:=|\theta_\star^{(-)}|\approx 1.2077065297.
$$
因此在“分歧半径”这一可计算的谱指纹下，负携带势的局部解析半径大于同步核输出位势，从而在相同的 $|\theta|$ 尺度上可获得更强的统一余项控制（Cauchy 估计中的 $r^{-n}$ 衰减更快）。
\end{remark}

\begin{remark}[高阶偶系数的振荡与最近复分歧角]\label{rem:pressure-even-derivative-oscillation}
令 $f(\theta):=P(\theta)-\theta/2-\log 3$，则 $f$ 为偶函数且在 $|\theta|<R_\theta$ 内解析。
由最近分歧点 $\theta_\star$ 的几何位置，可用 Darboux 型系数转移（或等价的 Cauchy 系数积分主贡献）解释高阶偶系数的符号/振荡：
其主导模式由 $\pm\theta_\star$ 及其共轭贡献叠加给出，故 $f^{(2k)}(0)$ 的符号可由
$\cos\bigl(2k\arg(\theta_\star)+\phi_0\bigr)$ 型项控制（$\phi_0$ 为分支点局部相位）。
在本核上 $\arg(\theta_\star)\approx 1.4017459236$。
\end{remark}
